



\usepackage[left=1in,right=1in]{geometry}
\usepackage{float}
\usepackage{enumitem}
\usepackage{fancyhdr}
\usepackage{pifont}
\usepackage{commath}
\usepackage{mathtools}
\usepackage[scr]{rsfso}
\usepackage[makeroom]{cancel}
\usepackage{bm}
\usepackage{fancybox,framed}
\usepackage{stackengine}
\usepackage{amssymb,amsmath}
\usepackage{amsthm}
%\usepackage{unicode-math}
\usepackage{esint}
%\usepackage{pxfonts}
\usepackage[overload]{empheq}
\usepackage{cases} 
\usepackage{accents} 
\usepackage[symbol]{footmisc}
\usepackage{cancel}
\usepackage{wrapfig}
\usepackage{graphicx}
\usepackage{subfigure}
\usepackage{xcolor}
\graphicspath{ {figures/} }
\usepackage[titletoc]{appendix}
\usepackage{tikz}
\usetikzlibrary{patterns,patterns.meta}
\usetikzlibrary {arrows.meta}
\usetikzlibrary{angles,quotes, calc}
\usetikzlibrary{decorations.pathreplacing,decorations.markings}
\usetikzlibrary {shapes.geometric}
\usepackage{tikz-3dplot}
\usepackage{pgfplots}
\pgfplotsset{width=10cm,compat=1.9}

\pgfdeclarepattern{
  name=hatch,
  parameters={\hatchsize,\hatchangle,\hatchlinewidth},
  bottom left={\pgfpoint{-.1pt}{-.1pt}},
  top right={\pgfpoint{\hatchsize+.1pt}{\hatchsize+.1pt}},
  tile size={\pgfpoint{\hatchsize}{\hatchsize}},
  tile transformation={\pgftransformrotate{\hatchangle}},
  code={
    \pgfsetlinewidth{\hatchlinewidth}
    \pgfpathmoveto{\pgfpoint{-.1pt}{-.1pt}}
    \pgfpathlineto{\pgfpoint{\hatchsize+.1pt}{\hatchsize+.1pt}}
    \pgfpathmoveto{\pgfpoint{-.1pt}{\hatchsize+.1pt}}
    \pgfpathlineto{\pgfpoint{\hatchsize+.1pt}{-.1pt}}
    \pgfusepath{stroke}
  }
}

\tikzset{
  hatch size/.store in=\hatchsize,
  hatch angle/.store in=\hatchangle,
  hatch line width/.store in=\hatchlinewidth,
  hatch size=5pt,
  hatch angle=0pt,
  hatch line width=.5pt,
}

\tikzset{
    tangent/.style={
        decoration={
            markings,% switch on markings
            mark=
                at position #1
                with
                {
                    \coordinate (tangent point-\pgfkeysvalueof{/pgf/decoration/mark info/sequence number}) at (0pt,0pt);
                    \coordinate (tangent unit vector-\pgfkeysvalueof{/pgf/decoration/mark info/sequence number}) at (1,0pt);
                    \coordinate (tangent orthogonal unit vector-\pgfkeysvalueof{/pgf/decoration/mark info/sequence number}) at (0pt,1);
                }
        },
        postaction=decorate
    },
    use tangent/.style={
        shift=(tangent point-#1),
        x=(tangent unit vector-#1),
        y=(tangent orthogonal unit vector-#1)
    },
    use tangent/.default=1
}
\tikzset{->-/.style={decoration={
  markings,
  mark=at position #1 with {\arrow{Stealth}}},postaction={decorate}}}


\tikzset{
    partial ellipse/.style args={#1:#2:#3}{
        insert path={+ (#1:#3) arc (#1:#2:#3)}
    }
}


\newcommand*\circled[1]{\tikz[baseline=(char.base)]{
            \node[shape=circle,draw,inner sep=2pt] (char) {#1};}}
\usepackage{pdfpages}

\newcommand{\q}[1]{``#1''}


\renewcommand{\thefootnote}{\fnsymbol{footnote}}






\theoremstyle{definition}
\newtheorem{definition}{Definition}[section]
\newtheorem*{example}{Example}
\newtheorem*{exercise}{Exercise}
\newtheorem*{solution}{Solution}
\newtheorem*{notation}{Notation}
\newtheorem*{note}{Note}
\newtheorem*{remark}{Remark}
\newtheorem*{discussion}{Discussion}
\newtheorem*{conclusion}{Conclusion}
%\newtheorem*{claimproof}{Proof of claim \theclaim}
%\newtheorem*{theoremproof}{Proof of theorem \thetheorem}






\theoremstyle{plain}
\newtheorem{theorem}{Theorem}[section]
\newtheorem{lemma}{Lemma}[section]
\newtheorem{claim}{Claim}[section]
\newtheorem{axiom}{Axiom}[section]



%\usepackage[thmmarks]{ntheorem}


%\theorembodyfont{\normalfont}
%\newtheorem{definition}{Definition}[section]
%\theoremstyle{plain}
%\theorembodyfont{\normalfont}
%\newtheorem{example}{Example}
%\theoremstyle{nonumberplain}
%\newtheorem{notation}{Notation}
%
%\theorembodyfont{\itshape}
%\theoremstyle{plain}
%\newtheorem{theorem}{Theorem}
%\newtheorem{lemma}[theorem]{Lemma}
%
%\newtheorem{claim}{Claim}[theorem]
%\theoremstyle{nonumberplain}
%\theoremsymbol{\ensuremath{\Box}}
%\newtheorem{proof}{Proof}
%\theoremsymbol{\ensuremath{\blacksquare}} 
%\newtheorem{claimproof}{Proof of claim \theclaim}

% argumented matrix
\makeatletter
\renewcommand*\env@matrix[1][*\c@MaxMatrixCols c]{%
  \hskip -\arraycolsep
  \let\@ifnextchar\new@ifnextchar
  \array{#1}}
\makeatother
%

% volume
\makeatletter
\DeclareRobustCommand{\volume}{\text{\volumedash}V}
\newcommand{\volumedash}{%
  \makebox[0pt][l]{%
    \ooalign{\hfil\hphantom{$\m@th V$}\hfil\cr\kern0.08em--\hfil\cr}%
  }%
}
%
\makeatother





\usepackage{hyperref}
\hypersetup{hidelinks,
	colorlinks=true,
	allcolors=black,
	pdfstartview=Fit,
	breaklinks=true
}

%------------Command----------------------------------------

\newcommand{\vc}[3]{\overset{#2}{\underset{#3}{#1}}}

%--------------tensor form---------------------
\stackMath
\newcommand\tenq[2][1]{%
 \def\useanchorwidth{T}%
  \ifnum#1>1%
    \stackunder[0pt]{\tenq[\numexpr#1-1\relax]{#2}}{\scriptscriptstyle\sim}%
  \else%
    \stackunder[1pt]{#2}{\scriptscriptstyle\sim}%
  \fi%
}

%-----------------------------------------

\DeclareMathOperator\erf{erf}
\DeclareMathOperator\erfc{erfc}

\DeclareMathOperator\Arg{Arg}
\DeclareMathOperator\Log{Log}

\DeclareMathOperator\arcsec{arcsec}
\DeclareMathOperator\arccot{arccot}
\DeclareMathOperator\arccsc{arccsc}
\DeclareMathOperator\tr{tr}

\def\Res{\mathop{\text{Res}}\limits}